\appendix

\section{附录}

\begin{frame}{附录 A：基准模型（只列核心结构）}
\small
\begin{block}{SN-PC$(\hat{\discretization})$}
\vspace{-0.25cm}
\begin{align*}
\min\;& \sum_{a\in \full{\flightArc}}\sum_{f\in\fleetSet} \cost_{af}\,\fleetAssignVar_{af}
- \sum_{m\in\marketSet}\sum_{p\in\psgTypeSet}\sum_{r\in\itinInMarket}\sum_{\ell\in\fareClassSet} \demand_{mp}\fare_{r\ell}\shareVar_{rp\ell}\\[0.15cm]
\text{s.t. } & \sum_{a\in \outArcSet{i,t}} \fleetAssignVar_{af} - \sum_{a\in \inArcSet{i,t}} \fleetAssignVar_{af} = 0, \quad \forall (i,t),f \\
& \sum_{a\in \holdingArcTc{\checkTime}} \fleetAssignVar_{af} + \sum_{a\in \flightArcTc{\checkTime}} \fleetAssignVar_{af} \le \fleetAvail_f, \quad \forall f \\
& \sum_{f\in\fleetSet}\sum_{a\in \arcDuringSepTime_{ijt}} \fleetAssignVar_{af} \le 1, \quad \forall (i,j),t \\
& \outAttr_{mp}\shareVar_{rp\ell} \le \attr_{rp\ell}\outsideShareVar_{mp}, \quad \forall m,p,r,\ell \\
& \left(\frac{\adjOutAttr_{mp}}{\outAttr_{mp}}\right) \outsideShareVar_{mp} +
\sum_{r,\ell} \left(\frac{\adjAttr_{rp\ell}}{\attr_{rp\ell}}\right)\shareVar_{rp\ell} = 1, \quad \forall m,p \\
& \sum_{m,p,\ell}\sum_{r\in \itinUsingArc{m}{a}} \demand_{mp}\shareVar_{rp\ell} \le \sum_{f\in\fleetSet}\fleetCap_f\fleetAssignVar_{af}, \quad \forall a\in\full{\flightArc}
\end{align*}
\end{block}
\end{frame}

\begin{frame}{附录 B：算法细节（按模块拆开）}
\small
\begin{columns}[T,totalwidth=\textwidth]
\column{0.33\textwidth}
\begin{block}{LBM}
\begin{itemize}
    \item 在粗时间离散化上建模；
    \item 允许时间弧聚合；
    \item 再对 itinerary 做 path-signature aggregation；
    \item 得到 valid lower bound。
\end{itemize}
\end{block}

\column{0.33\textwidth}
\begin{block}{UBM}
\begin{itemize}
    \item 在 full network 上修复；
    \item 锁定 segment--fleet frequency；
    \item 用 itinerary cuts 抑制收入高估；
    \item 给出 feasible incumbent。
\end{itemize}
\end{block}

\column{0.33\textwidth}
\begin{block}{Refinement}
\begin{itemize}
    \item 检查 chronological violations；
    \item 检查 separation infeasibility；
    \item 检查 phantom revenue；
    \item 针对性加入时间点，继续迭代。
\end{itemize}
\end{block}
\end{columns}

\vspace{0.2cm}
\begin{alertblock}{一句话总结}
\centering
\textbf{用 tighter LB 缩小搜索空间，用 targeted refinement 保证 exactness。}
\end{alertblock}
\end{frame}

\begin{frame}{附录 C：行程聚合为什么能保持下界？}
\begin{columns}[T,totalwidth=\textwidth]
\column{0.42\textwidth}
\begin{block}{直观理解}
\small
\begin{itemize}
    \item 一个 super-itinerary 对应多个原始 itinerary；
    \item 原始 itinerary 之间可能因最小间隔约束而冲突；
    \item 我们不是把 attraction 全部直接相加；
    \item 而是在合法组合上取 \hl{optimistic but valid} 的上界。
\end{itemize}
\end{block}

\vspace{0.12cm}
\begin{equation*}
\attrSuper_{R p\ell} =
\max_{\itinCombSet:\, |\itinCombSet|\le \combLimit_R}
\sum_{\mathscr{O}\in\itinCombSet}\sum_{r\in\mathscr{O}} \attr_{r p\ell}
\end{equation*}

\vspace{0.05cm}
\small
其中 $\combLimit_R$ 由聚合时间窗内可物理共存的服务数量控制。

\column{0.56\textwidth}
\begin{center}
    \includegraphics[width=\linewidth]{../figure/illustrate_valid_combinations/valid_combinations.pdf}
\end{center}
\end{columns}
\end{frame}
